\section{Introduction}\label{sec:introduction}

The proliferation of autonomous multi-agent AI systems has produced a structural crisis in computational accountability. As agent networks grow in depth, recursion, and asynchronous coordination, the human principal who authorizes a system's operation finds herself increasingly distant from — and structurally unreachable by — the actions that the system performs in her name. Delegation chains composed entirely of machine actors terminate in outcomes for which no human can be held architecturally responsible. We call this the \emph{accountability orphan} problem: actions performed by autonomous systems that are causally downstream of human authorization but operationally severed from it.

The consequences of the accountability orphan problem are not merely philosophical. When a multi-agent system produces a harmful outcome — a misallocated resource, a safety-critical misconfiguration, an irreversible physical action — and that outcome is mediated by multiple layers of autonomous decision-making, the standard remediation pathway collapses. No individual machine actor can be held responsible, because each was operating within its provisioned scope. No human actor can be held responsible, because no human approved or could foresee the specific outcome. The accountability gap is not a governance failure; it is an architectural one. It arises because the system was designed without a mandatory mechanism that guarantees human reachability under all exception conditions, including those that were not anticipated at design time.

Existing approaches to this problem fall into three categories, each of which we argue is insufficient on its own. The first is post-hoc audit: maintaining detailed logs of agent decisions for human review after execution. Post-hoc audit is necessary but not sufficient. A historical record that cannot enforce its own conclusions — that has no pathway to actually stop a running system or reverse a harmful trajectory — is accountability theater. The second is static permission scoping: defining hard boundaries on what each agent may do, enforced by capability lists or role-based access controls. Static scoping fails because it assumes that the space of possible exception conditions is enumerable at design time, an assumption that does not hold for systems operating in open environments with learned behavioral models. The third is human-in-the-loop (HITL) review, in which humans intervene at designated checkpoints in the agent workflow. HITL is directionally correct but typically implemented as a discretionary convention rather than an architectural compulsion. A human who chooses not to intervene, or who is unavailable, or who is herself operating under time pressure, becomes a薄弱 link — a single point of failure that the system cannot bypass.

The core insight of this paper is that accountability must have an \emph{architectural fallback}. It is not enough for a system to provide pathways for human intervention; the system must guarantee, by construction, that every unresolved exception condition — regardless of its novelty, cause, or origin — reaches a designated human principal. This guarantee cannot be implemented as a policy or a convention. It must be embedded in the execution architecture as a deterministic, non-negotiable property: upon any condition that a node cannot resolve with certainty, the node must transition to a state in which no further autonomous harm is possible and the exception must propagate to a human accountability anchor. No machine actor along that propagation path may absorb, resolve, or suppress the exception. This is not a preference; it is an architectural invariant.

We call this fallback mechanism the \textbf{Bailout Protocol}. It is implemented as a deterministic finite state machine embedded in every node of a \bxthree{}-architected multi-agent system, operating across the \purpose{}, \bounds{}, and \fact{} layers. The \purpose{} layer carries final authority and receives all escalated exceptions; the \bounds{} layer evaluates whether proposed actions fall within a node's provisioned scope and initiates bailout when they do not; the \fact{} layer provides deterministic enforcement, cryptographic sealing of the forensic record, and the immutable transition relation governing the protocol's state machine. The protocol is triggered not by any specific anticipated failure mode but by a single, general condition: the node's certainty threshold for safe autonomous action has been violated. This generality is essential. Anticipatory trigger conditions are inherently incomplete; a bailout mechanism that only handles enumerated failure modes is structurally vulnerable to the unanticipated ones that matter most.

The Bailout Protocol contributes to the literature in three ways. First, we formalize the \emph{accountability orphan} problem as an architectural gap in multi-agent AI systems, distinguishing it from governance, audit, and access-control approaches that address different dimensions of the same underlying risk. Second, we specify the Bailout Protocol as a deterministic state machine with formally defined trigger conditions, transition relations, escalation algorithms, and timeout guarantees, providing a machine-verifiable implementation substrate. Third, we demonstrate that accountability guarantees can be made architecturally non-negotiable — that human reachability can survive delegation depth, asynchronous operation, network partition, and recursive node spawning — by exploiting the immutable parent-pointer structure and human accountability anchor reference embedded in each node's worksheet.

The paper is organized as follows. Section~\ref{sec:background} reviews the three lineages of research that inform the Bailout Protocol: accountability in multi-agent AI systems, human-in-the-loop design, and fail-safe design in safety-critical systems. It establishes the \bxthree{} three-layer model as the architectural substrate against which the protocol operates. Section~\ref{sec:state-machine} provides the complete formal specification of the protocol as a deterministic finite state machine, including the state space, input alphabet, transition relation, bypass mechanism, and timeout handling. Section~\ref{sec:trigger} defines the three trigger conditions — Capability Boundary, Safety Envelope Violation, and Accountability Boundary Violation — and the certainty-threshold mechanism that generalizes them. Section~\ref{sec:cascading} describes the Cascading Triggers system for propagating exceptions across delegation chains, including the bypass mechanism and the conditions under which propagation routes to an alternate functional pathway. Section~\ref{sec:forensic} specifies the Forensic Ledger interface and the append-only sealing protocol maintained by the \fact{} layer. Section~\ref{sec:implementation} presents the implementation architecture and evaluation results across three deployed systems: a precision agriculture platform, an enterprise workforce orchestration system, and a sovereign AI deployment. Section~\ref{sec:related} reviews related work. Section~\ref{sec:conclusion} concludes with a discussion of limitations and directions for future research.
